\documentclass[11pt]{article}
\usepackage{float}
\usepackage[most]{tcolorbox}
\usepackage{amsmath,amssymb} 
\usepackage{fourier} 
\usepackage{xcolor}
\definecolor{darkblue}{RGB}{128, 0, 32} 
\newcommand{\D}[2]{\frac{\mathrm{D}#1}{\mathrm{D}#2}}

\newcommand{\eqset}[1]{
\begin{spreadlines}{8pt}
\begin{empheq}[left =\empheqlparen]{equation}
\begin{aligned}
#1
\end{aligned}
\end{empheq}
\end{spreadlines}
}

\newcommand{\headerfont}{\fontfamily{ppl}\selectfont} 

\usepackage{physics}
\let\braket\undefined
\newcommand{\braket}[2]{\left\langle #1 \middle| #2 \right\rangle}
\usepackage{siunitx}
\AtBeginDocument{\RenewCommandCopy\qty\SI}
\usepackage{slashed}
\usepackage{geometry}
\geometry{top = 2.5cm, right = 2cm, left = 2cm, bottom = 2.5cm} 

\usepackage{hyperref}
\hypersetup{
    colorlinks=true, 
    linkcolor=darkblue, 
    urlcolor=darkblue,
    citecolor=magenta
}

\usepackage[explicit]{titlesec}
\usepackage{tabto}
\renewcommand{\thesubsection}{\alph{subsection}} 

\titleformat{\section}[block]
  {\headerfont\color{darkblue}\Large} 
  {} 
  {0pt} 
  {\fcolorbox{darkblue}{white}{\parbox{\dimexpr\textwidth-2\fboxsep-2\fboxrule\relax}{\thesection\hspace{1em}#1}}}
\usepackage{accents}
\renewcommand{\vec}[1]{\underaccent{\widetilde}{#1}}
\titleformat{\subsection}{\headerfont\color{darkblue}\large}{\thesubsection}{1em}{#1}

\usepackage{fancyhdr}
\pagestyle{fancy}
\fancyhf{}
\renewcommand{\headrulewidth}{0pt}
\renewcommand{\vec}[1]{\underline{\textbf{#1}}}
\renewcommand{\sectionmark}[1]{\markboth{#1}{}} 

\newcommand{\entrytitle}{Liquids and Glasses}
\fancyhead[C]{%
    \colorbox{darkblue}{%
        \parbox{\dimexpr\textwidth-2\fboxsep\relax}{%
            \headerfont\color{white} 
            \entrytitle \hfill \leftmark 
        }%
    }%
}
\fancyfoot[C]{\thepage}

\newcommand{\eq}[1]{\begin{equation} #1 \end{equation}}
\newtcolorbox{customBox}{
    colback = white, 
    colframe = black, 
    boxrule = 0pt, 
    leftrule = 6pt,
    sharp corners
}
\newcommand{\avg}[1]{\langle #1 \rangle }

\newcommand{\defn}[1]{\begin{customBox}#1\end{customBox}}
\newcommand{\thm}[1]{
    \begin{customBox}
        \ttfamily
        #1
    \end{customBox}
}


\newcommand{\dbar}{d\hspace*{-0.08em}\bar{}\hspace*{0.1em}}

\newcommand{\ex}[1]{\begin{customBox}\textit{Example:} #1\end{customBox}}
\DeclareRobustCommand{\rchi}{{\mathpalette\irchi\relax}}
\newcommand{\irchi}[2]{\raisebox{\depth}{$#1\chi$}}

\title{\headerfont\color{darkblue}\Huge \textbf{Title}}
\author{Author}
\date{\headerfont\today}

\begin{document}

\maketitle
\thispagestyle{empty} 
\vfill 

\begin{tcolorbox}[
    colback=gray!10,  
    colframe=gray!10, 
    boxrule=0pt,     
    sharp corners,    
    left=15pt, right=15pt, top=15pt, bottom=15pt 
]
{\headerfont\color{darkblue}\large\textbf{About}}\par\vspace{0.5em}
Notes on [topic].
Note that vectors corresponding to single particles are bolded (for example, \textbf{r}) while vectors over the entire system (i.e. $N\times d$ dimensional vectors) are bolded and underlined: $\vec{r}$.
\end{tcolorbox}
\newpage

\tableofcontents
\thispagestyle{empty}

\newpage
\setcounter{page}{1} 

\section{Section}








\bibliographystyle{plain}
\bibliography{biblio}
\appendix
\section{Appendix section}





\end{document}























































































































































































































































































































































































































































































































































































































































































































































































































































































































































































































































































`

